\begin{abstract}
An LLM asked to forget specific information should preserve unrelated
knowledge and capabilities. Yet approximate unlearning can cause concentrated
collateral damage that aggregate scores hide. We ask whether susceptible
retained behaviors can be identified before a target unlearning run and
selectively protected without modifying the parent unlearner.

We combine two model-derived signals that require no candidate-side
backpropagation during target profiling. \emph{Loss Susceptibility} is a
gradient-defined quantity: the maximum squared first-order loss response over
unit-norm perturbations within a declared parameter block, without conditioning
on the later update direction. Its pool-shared forward-only estimator uses
shared symmetric perturbations.
\emph{Representation Proximity} estimates the behavior's exposure to the
current forget request in hidden-state space. During target profiling, both
signal estimates use forward evaluations only: perturbations are shared across
the candidate pool and hidden states are reused, so no candidate-specific
gradient is materialized. A scalar mixture weight is calibrated on
target-request-disjoint development trajectories, and the resulting score is
sealed before the target run.

Together, these signals turn collateral damage from a failure observed after
unlearning into a risk that can be anticipated and managed. The forward-only
profile requires no realized target update, so it can screen retained
behaviors before the target trajectory. When exhaustive auditing or repair is
impractical, it narrows a large candidate universe to those most worth
monitoring or protecting. Under a fixed quota and repair budget, protection
becomes an allocation question: where should limited capacity be spent? We
evaluate held-out damage prediction, the forward-only estimator's fidelity to
exact Loss Susceptibility on each setting's declared fidelity support, and
repair selection under one feasibility contract. \resph{Report the resolved
prediction, fidelity, protection, and failure-boundary results.} This connects
early warning to selective protection across models and unlearning objectives
without modifying the parent method or claiming a deletion certificate.

\end{abstract}
